\documentclass[11pt]{article}
\usepackage[margin=1in]{geometry}
\usepackage{amsmath,amssymb}
\usepackage{booktabs}
\usepackage{xcolor}
\usepackage[colorlinks=true,linkcolor=black,urlcolor=blue]{hyperref}
\usepackage{parskip}
\newcommand{\xor}{\oplus}
\newcommand{\code}[1]{\texttt{#1}}
\title{\textbf{The single-carrier product-share gadgetizer}\\[2pt]
\large Construction, bookkeeping, and the statistical measurements}
\author{}
\date{2026-07-26}
\begin{document}
\maketitle

\noindent\emph{Assumes} \code{docs/PRODUCT\_SHARE\_ENCODING}, \code{docs/BAND\_HARDENING}
\emph{and} \code{docs/NONLINEAR\_GADGETIZATION}. \emph{This document describes the
construction as it now stands, including the single-carrier decode, and
presents the measurements that support it --- in particular the tradeoff
between the two shipped mask plans.}

\section{What the gadgetizer builds, and against whom}

The gadgetizer rewrites a circuit $A$ (in production, a sliced sandwich of
width $2n$ wrapping the source circuit $C$) into a much larger equivalent $G$
that computes the same function on its low wires when every other wire enters
at $0$. The security goal is that $G$ leak nothing about $A$ beyond its
input/output behaviour, and the specific leak this line of work attacks is the
\textbf{computational-progress diagonal}: the ability to tell \emph{where in
the original computation} a point of $G$ sits.

Two execution-based instruments measure it, and they answer different
questions. \code{hmap\_affine} asks whether $C$'s state at prefix $i$ is an
\emph{exact} GF(2) function of degree $\le d$ of $G$'s wires at prefix $j$,
scoring every inexact cell at $0.5$; a low-error valley advancing with $i$ is
the diagonal. \code{hmap\_stat} asks the complementary question --- how well
can a cheap predictor \emph{guess} the bit --- and is sensitive to bias, which
the exact-span measure is structurally blind to. Both are read by the ridge
measure, and on plates with unencoded ports they must be read by
\textbf{interior row structure}, never by the global $\rho$
(\S\ref{sec:method}).

\section{The encoding}

Each logical value $v_i$ is carried as

\begin{equation}
v_i \;=\; \underbrace{w_{u(i)} \;[\,\xor\; w_{u'(i)}\,]}_{\text{one or two linear carriers}}
\;\xor\; \underbrace{\bigoplus_{j} \prod_{l} \bigl(\mathrm{loc}[b_{jl}] \xor a_{jl}\bigr)}_{k_{\mathrm{total}}\ \text{product mask terms}}
\;\xor\; c_i ,
\label{eq:decode}
\end{equation}

where each mask term is a conjunction of $\deg$ band-variable literals with
compile-time offsets $a_{jl}$, and $c_i$ is a compile-time ledger constant. The
$a$ and $c$ bits never appear on a wire.

\subsection{Why a linear term is forced, and why one is enough}

\textbf{Balance obstruction.} No exact reversible gadget can conditionally flip
a value whose decode has unequal representation-class sizes, with any ancillae,
garbage or re-encoding: exactness forces a bijection of
$\{v_x=\alpha, v_t=\beta\}$ onto $\{v_x=\alpha, v_t=\beta\xor\alpha\}$, and
counting representations gives $R'(0)/R'(1)=R(0)/R(1)$ from the $\alpha=0$
equations and $R'(0)/R'(1)=R(1)/R(0)$ from the $\alpha=1$ equations, jointly
forcing $R(0)=R(1)$. A pure two-wire product decode is $3{:}1$, hence
ungateable. Since every \emph{balanced} two-wire decode is affine, the
nonlinearity must come from wires beyond the value's own pair.

\textbf{But that argument only excludes having \emph{no} linear part.} Write
the decode as $D = c \xor g(\text{sources})$ and fix the sources arbitrarily:
flipping $c$ flips $D$, so the two classes are exactly equal for \emph{any}
$g$. \textbf{One linear term discharges the obstruction completely.}

\textbf{And the second linear term is free to the adversary.} An affine
adversary obtains every degree-1 term at no cost by XORing it into its
predictor, so the second carrier never enters the piling-up product, which runs
over the \emph{nonlinear} terms alone. Writing a plan as the multiset of term
degrees, the shipped \code{[3,3]} decode is really $[1,1,3,3]$, and the second
$1$ buys nothing statistically.

\subsection{Plan notation and the piling-up value}

Under jointly uniform sources with variable-disjoint terms, the naive predictor
(the XOR of the linear carriers) agrees with the value with probability

\begin{equation}
\tfrac12 + \tfrac12 \prod_{j \,:\, d_j > 1} \bigl(1 - 2^{\,1-d_j}\bigr),
\label{eq:piling}
\end{equation}

the product running over the nonlinear terms only. This gives $0.781$ for
$[1,1,3,3]$ and $[1,3,3]$ alike, $0.641$ for $[1,2,3,3]$, and $0.594$ for
$[1,2,2,3]$. Section~\ref{sec:limits} records the hypothesis under which
\eqref{eq:piling} is a theorem rather than a model, and why it is not yet met.

\section{W: the sharing bookend and the ports}

\subsection{Two-carrier build}

The input port emits, in order: a randomising bookend, the band fill, the $W_i$
sharing ramp, and the mask injection ramp.

\textbf{Bookend.} \code{rand\_z\_gates} emits g57 gates whose \emph{targets} are
drawn round-robin from the aux half $[n,2n)$ --- each aux wire is a target
exactly once per $n$-gate round --- while both controls are drawn uniformly
from \emph{all} $2n$ wires. No data wire is ever a target, so the low $n$ wires
still hold $x$ when the band fill and the ramp run. Its size is
$\max(2n\lfloor \ln n\rfloor, 64)$ gates (note the floor: at $n{=}128$ this is
$1024$, not $1243$).

\textbf{$W_i$ ramp.} For each value, \code{emit\_w\_i\_cnot} emits seven width-1
CNOTs realising the GF(2)-linear map $(a,b,c,d)\mapsto(c,\,d,\,a\xor b,\,b)$ on
(data, aux, share, pad). The effect is that the wire chosen as \code{share}
receives $v \xor z$ and the wire chosen as \code{pad} receives $z$, so
$\text{share}\xor\text{pad}=v$; simultaneously the old contents of those two
wires are relocated onto the vacated data and aux wires. \code{share} and
\code{pad} are drawn by rejection uniformly from all $2n$ carriers, never from
the band, which is exactly why relocation bookkeeping is needed: a draw may
land on a wire currently holding another value's material. That bookkeeping is
the \code{Slot} enum (\code{Data}/\code{Aux}/\code{Pair}/\code{Output}) plus
\code{reloc}, which rewrites \code{dloc}, \code{aloc} or the matching half of
\code{pairs}.

\textbf{Decode bookend.} The output port runs $W_i^{-1}$ (the same seven CNOTs
reversed, each transvection being self-inverse), borrowing a carrier when
neither the share nor the pad already sits on the value's home wire.

\subsection{Single-carrier build}

There is \textbf{no $W_i$ ramp and no \code{rand\_z} bookend at all}. The band
fill is the first thing emitted; \code{pairs[v]} is initialised to the
degenerate tuple $(v,v)$, so value $v$ simply sits on wire $v$; and
\code{inject\_all} --- one conjunction per planned mask term, $n\,k_{\mathrm{total}}$
in total --- \emph{is} the entire ramp. The value is masked in place. Exit is
the mask strip followed by a routing permutation (\S\ref{sec:rg}).

This makes the single-carrier construction structurally \emph{simpler} than the
paired one, and it is worth being explicit that all port-side hiding then comes
from the mask injection rather than from a randomised sharing.

\subsection{The band fill}

The \textbf{linear} fill gives each band wire its input junk XOR
$\langle\alpha,x\rangle$, with $\alpha$ uniform on $\{0,1\}^n$ conditioned on
Hamming weight $\ge 2$ (the subset is resampled until it has at least two
members). The \textbf{nonlinear cascade} fill instead emits, per band wire in
list order,
\[
\text{band}_j \;=\; \text{junk} \;\xor\; x_{p_j} \;\xor\; (\text{small linear part}) \;\xor\; \bigoplus^{M} (\text{two-source products}),
\]
where the product sources are drawn from data wires and from \emph{earlier}
band wires --- ``earlier'' meaning earlier \emph{position in the wire list}
passed in, not lower wire index, since after rolling that list is an arbitrary
set. A transitive-support map is maintained per filled wire so that a cascade
reference which would drag in this wire's own pivot is filtered out. The
cascade is what reaches high input degree with nothing wider than a two-control
gate; a flat high-degree monomial would fire on a $2^{-d}$ fraction of inputs
and behave like its linear part.

\textbf{Balance.} The pivot appears exactly once, linearly, and is excluded
from the linear pool, from the data branch of the product draw, and
transitively from the band branch. Hence each band bit is exactly balanced over
uniform $x$ for \emph{any} choice of the nonlinear part. This is a
\emph{marginal} guarantee, one wire at a time, and \S\ref{sec:limits} explains
why that is weaker than the mask statistics actually need.

\textbf{Mirror fill.} The same fill is emitted again after the output bookend.
In the paired build its target set is \emph{every non-output wire} --- the aux
half plus the band, not the band alone --- precisely because the emitted target
set is part of the gate list, and filling only the band's final wires would
publish where the rolls left the band.

\subsection{The zero-slice preblock}

The preblock makes the gadget compute $A$ \emph{only} on the all-zero slice.
Every gate is a positive conjunction whose target is a data wire and which
reads exactly one slice wire, so every gate is individually dead when the slice
is zero; slice controls are assigned balanced (each non-data wire is read at
least once) and then shuffled. A third of the gates are plain CNOTs, and the
remainder split between one and two data controls --- the two-data-control
shape is what keeps the disturbance from being affine in $x$.

The acceptance predicate is that \emph{no nonzero slice acts as the identity on
the data}: for every nonzero slice there \emph{exists} at least one disturbed
input. It is checked exhaustively when $n + \text{slice} \le 20$ and by
sampling otherwise, redrawing until satisfied.

Dimensions differ by build: the paired slice is $n + \text{band}$ (aux half
plus band, total $2n+\text{band}$), the single-carrier slice is the
\textbf{band alone} (total $n+\text{band}$).

\section{CG: the share-native fold}

A source gate is $\code{target} \mathrel{\xor}= \code{comp} \xor \prod_i
\mathrm{lit}_i$. The fold substitutes each operand's decode
\eqref{eq:decode} by its \emph{atom list}, walks the cartesian product, and
emits each cross term as one conjunction fragment.

\textbf{Atoms.} For each control $(w,\text{positive})$: the linear atoms
(\emph{two} singleton atoms $[(c_0,\top)]$ and $[(c_1,\top)]$ in the paired
build; exactly \emph{one} under a single carrier, since emitting it twice would
cancel); one atom per live mask slot, resolved to $\deg$ physical literals; and
a \emph{constant} atom (the empty vector) exactly when
$\code{consts}[w] \xor \lnot\text{positive}$ --- the ledger constant of $w$
parity-merged with the control's own polarity. This is the only place control
polarity enters.

\textbf{Fragments.} Each cross term is the concatenation of one atom per
operand. A term that is entirely empty is not emitted: it flips the target's
ledger constant instead (this is also how an arity-0 source, an X/NOT gate,
becomes pure bookkeeping). A term with two polarities of one wire is
identically zero and is dropped. Every surviving fragment is written to one
carrier of the target --- a coin flip in the paired build, the unnamed carrier
under distributed sourcing, the single carrier otherwise.

\textbf{Counts.} The fragment count per source gate is the product over
operands of their atom counts:
\[
(2 + k_{\mathrm{total}})^{\text{arity}} \quad\text{two-carrier},
\qquad
(1 + k_{\mathrm{total}})^{\text{arity}} \quad\text{single-carrier},
\]
at equal mask strength. Fragment \emph{width} is set by which atoms are picked
(1 for a carrier atom, $\deg$ for a mask atom, 0 for the constant atom), so it
reaches $\text{arity} \cdot \max\deg$.

\textbf{No bare X.} Every constant is ledger-only: the gate complement, the
all-constant cross term, the residual constant of a narrow realisation, and the
slot injection/strip parities. The accumulated constant is discharged once, in
the final strip, as the two-gate pair $\lnot u$ then $u$ over a random helper
carrier --- never as an X gate.

\textbf{Emission order.} Fragments are materialised, \emph{shuffled}, and only
then assigned targets. The odometer order is a static per-gate progress clock:
consecutive fragments share atom prefixes, and the fold blocks would be legible
with no execution at all.

\textbf{Gate-local non-completeness.} One atom per operand means a fragment
reads at most one carrier of each operand and never a full decode; mask atoms
resolve to band wires, which the roll keeps disjoint from carriers; and the
target is never among its own controls. Under distributed sourcing this is no
longer structural and is asserted at emission time instead
(\code{debug\_check\_fragment}, a debug-build guarantee).

\textbf{Narrow mode.} \code{--prod-max-width} ladders every emission down to a
width cap using a Barenco-style double sweep over \emph{borrowed dirty}
carriers: each borrow is visited an even number of times, so its unknown
content cancels, which costs no extra wires and preserves the arbitrary-junk
contract. Rungs must be exact width-2 conjunctions --- a g57's complement on a
borrowed wire does not cancel. Narrow mode is built and measured but
\textbf{not recommended}: it costs $6.4\times$ the gates, and the premise
(that wide fragments would be inert fossils) was falsified --- wide mask
fragments are the \emph{most} mobile gates in the circuit.

\section{RG: re-randomisation}
\label{sec:rg}

\subsection{The paired build}

The production path draws uniformly from three nonlinear g57 networks, once per
inter-source-gate gap by default.

\begin{center}
\begin{tabular}{llp{0.52\textwidth}}
\toprule
network & writes & effect \\
\midrule
RG1 (deg 3) & all four wires of two pairs & swaps which carrier \emph{pair} holds which value ($\code{pairs.swap}(i,j)$); the value decodes identically, off the other pair's wires \\
RG2 (deg 2) & all four wires & re-pairs: value $i$ moves to (pad$_i$, share$_j$) and $j$ to (share$_i$, pad$_j$) \\
RG3 & both carriers of one value & XORs the identical nonlinear bit into both carriers, so the carrier XOR --- the value --- is preserved \\
\bottomrule
\end{tabular}
\end{center}

All three preserve every logical value, and all three leave the masks invariant
because slots name band \emph{variables} rather than wires (\S\ref{sec:ledger}).

\textbf{A deliberate trade worth stating plainly.} The \emph{linear} CNOT
variants of these networks have gate-local non-completeness (no gate reads both
carriers of a value; a bounded exhaustive search shows six CNOTs is minimal for
RG1 with that property). The \emph{nonlinear} networks that actually ship
\textbf{do not}: RG1's every gate reads both carriers of one value, and RG3's
two random controls may be the two carriers of a third value. Non-completeness
was given up for low-degree opacity. It would be exactly backwards to describe
the production RGs as non-complete.

\subsection{The band roll}

A roll picks a band variable and a uniformly random partner wire and swaps
their contents with three transvections; the bookkeeping then re-points, so
either two band variables trade wires, or a band variable and a carrier
\emph{trade roles}. Nothing about the encoding changes --- the band variable
keeps its value, the carrier keeps its value --- only where the band lives.
Carriers and band wires remain a partition, so a fold or inject target is never
also a mask literal.

Each transvection is drawn from a mixed vocabulary: either a plain CNOT or the
pair $t \mathrel{\xor}= s \wedge u$, $t \mathrel{\xor}= s \wedge \lnot u$ over a
random helper, which sums to exactly $s$ and adds no constant. A roll is
therefore 3--6 gates, not 3. The reason is a gate-\emph{shape} census: before
this change, nothing but rolls ever wrote a wire while it held a band variable,
so those wires were written almost exclusively by width-1 CNOTs and remained
separable even once the write \emph{counts} matched.

\subsection{Single-carrier: R1 and re-sourcing}

RG2 and RG3 have \textbf{no single-carrier analogue}. There is no pair to
re-pair; and XORing into a lone carrier changes the value, which could only be
compensated by a term equal to the injected bit --- that is, by re-creating the
second carrier. \emph{The two-carrier structure is the refresh mechanism.}

What remains covers both roles:

\begin{itemize}
\item \textbf{Position} --- \code{emit\_value\_relocation} picks two distinct
\emph{values}, swaps their wires with a three-CNOT swap, and exchanges the
bookkeeping. It deliberately makes no carrier-space assumption, because a roll
can leave a value outside the carrier range.
\item \textbf{Representation} --- mask re-sourcing (\code{--prod-rsrc}) injects
a fresh term and strips the old one, which genuinely changes the carrier's bit.
This makes \code{--prod-rsrc} $\ge 1$ load-bearing in single-carrier mode; it is
documented but not asserted.
\end{itemize}

Because relocations and rolls permute values across the \emph{whole} wire
space, the exit emits a full permutation routing --- cycle resolution over all
wires, not a carrier-space permutation --- so that wires $[0,n)$ hold the values
and the band range holds junk, which is what makes the mirror fill safe.

\section{The ledger: tracking the high-degree share elements}
\label{sec:ledger}

This is the bookkeeping that makes a moving band possible.

\textbf{A slot names variables, not wires.} A \code{ProdSlot} is a vector of
factors $(b, a)$ where $b$ is a \emph{band-variable id}. The single place a slot
becomes physical is
\[
\code{lits}(\mathrm{loc}) : (b,a) \;\mapsto\; (\mathrm{loc}[b],\ \lnot a),
\]
evaluated at emission time. \code{loc} maps variable id to current wire; it
starts as the contiguous home range and changes \emph{only} in \code{roll}. This
indirection is exactly what makes a relocation free: the strip cancels the same
product the inject added, because both resolve through the same map. Under
distributed sourcing \code{loc} is empty and \code{lits} treats the factor id as
a physical wire (the identity map) --- so it is wrong to say literals are
\emph{always} resolved through \code{loc}.

\textbf{Per-value state.} A shared \code{plan} vector (the multiset of degrees,
$k$ copies of $\deg$ then $k_{hi}$ of $\deg_{hi}$), a per-value list of live
slots, and a per-value constant. Slots are drawn against a \emph{global} dedup
set keyed on the sorted variable tuple \emph{together with its offset pattern},
which is why the sizing bound counts $\binom{b}{d}2^{d}$ rather than
$\binom{b}{d}$.

\textbf{Lifecycle.} Draw (uniform distinct variables, sorted for canonical
form, independent uniform offsets, rejected if already live, hard panic on
exhaustion) $\to$ inject $\to$ optional re-source $\to$ strip. Injection and
strip are the \emph{same} routine, which is what makes a slot self-inverse up to
its returned constant parity. Re-sourcing always injects the replacement
\emph{before} stripping the old term, so a value is never momentarily bare. The
final strip walks slots in plan order (base terms first), then discharges any
residual constant.

\textbf{Invariants.} Carriers and band wires partition the wire space at all
times; a roll exchanges the two roles rather than duplicating them. No slot
names a carrier of its own value. Every emission's residual parity is folded
into the value's ledger constant. Slot identity is invariant under rolls.

\textbf{Distributed sourcing (off by default).} When mask factors live on
ordinary carriers instead of a band, three further mechanisms appear: a
reference count per wire; \emph{invariant S} --- at most one carrier of any
value is named by live slots, which guarantees every value always has a
writable carrier and prevents any fragment from seeing both carriers of a third
value; one factor per owning value within a slot; and a write barrier that
re-sources every slot naming a wire before anything writes it, forbidding the
whole released set to the replacement draw. A 64-lane simulation of the emitted
prefix lets the draw reject factor sets that are degenerate as \emph{bits}
(equal or complementary wires, identically-zero products) and replacements that
duplicate a term the value already carries. Section~\ref{sec:dist} reports why
this mode is not recommended.

\section{Methodology: reading the plates}
\label{sec:method}

Three conventions are load-bearing, and each was learned by getting it wrong.

\textbf{Read interior rows, not $\rho$.} Both axes have unencoded ends: cell
$(0,0)$ compares $C$'s input state against $G$'s input wires and reads
perfectly in every build, encoded or not. On a plate with port plateaus the
global $\rho$ is that two-plateau shape ranked against noise in a dead middle,
and its value is close to a coin flip. Report the count of interior rows with
any prominence.

\textbf{A capability control is mandatory for degree-$d$ claims.} A
degree-2 adversary that cannot break a degree-2-masked build is too weak to
conclude anything from. Two attempts at the degree-2 comparison in
\S\ref{sec:tradeoff} were void for exactly this reason --- once because the
adversary was restricted to 40 of 256 band wires (it then holds a given
degree-3 mask's factors about $0.4\%$ of the time), and once because rolls had
moved the band out from under a fixed wire list. Both reported ``dead'' for
every build, which would have been a false all-clear.

\textbf{Match the sampling.} The ridge statistic is a per-row maximum over
columns, so plates with different column counts are not comparable; choose the
stride per build to equalise them.

\section{Measurements}

All builds below share the source circuit and, at $n{=}128$, a byte-identical
sandwich; only the encoding differs.

\subsection{The decode grid at $n{=}16$}

Ridge is degree-1 interior median prominence; agreement is \code{hmap\_stat}'s
interior median (a lower bound: the predictor family stops at two wires).

\begin{center}
\begin{tabular}{lrrrrr}
\toprule
decode & gates & wires & ridge & agreement (peak) & theory \\
\midrule
plain (no masks) & --- & 32 & 0.3750 & --- & 1.000 \\
$[1,1,3,3]$ two-carrier & 4\,848 & 44 & 0.0312 & 0.720 (0.877) & 0.781 \\
$[1,3,3]$ single & 2\,866 & \textbf{28} & 0.0312 & 0.757 & 0.781 \\
$[1,2,3,3]$ single & 4\,063 & 30 & 0.0312 & \textbf{0.655} (0.680) & 0.641 \\
$[1,2,2,3]$ single & 4\,115 & 30 & 0.0312 & \textbf{0.618} (0.643) & 0.594 \\
\bottomrule
\end{tabular}
\end{center}

The ridge does not move: every encoded build is identically dead against the
plain control's $0.3750$. The measured agreements track \eqref{eq:piling} to
within a few points, which is itself a check on the arithmetic.

Note $[1,3,3]$ measuring \emph{worse} than $[1,1,3,3]$ at identical theory
($0.757$ vs $0.720$). That is the probe-locality cost surfacing: with one
carrier the single wire \emph{is} the predictor, so no pair search is needed
and the measure lands closer to the true value. One degree-2 term more than
repays it.

\subsection{Production, $n{=}128$ sliced sandwich}
\label{sec:prod}

\begin{center}
\begin{tabular}{lrrrrr}
\toprule
build & gates & wires & ridge (max) & $\rho$ & agreement (peak) \\
\midrule
$[3,3]$ two-carrier + roll & 247\,326 & 558 & 0.0000 (0.0547) & 0.039 & 0.719 (0.749) \\
$[1,2,3,3]$ single, band 56 & \textbf{223\,379} & \textbf{312} & 0.0000 (0.0547) & 0.016 & 0.632 (0.660) \\
$[1,2,3,3]$ single, band 256 & 229\,712 & 512 & 0.0000 (\textbf{0.0430}) & 0.016 & 0.618 (0.652) \\
$[1,2,2,3]$ single, band 256 & 230\,049 & 512 & 0.0000 (0.0430) & 0.015 & \textbf{0.583 (0.600)} \\
\bottomrule
\end{tabular}
\end{center}

All four are dead at degree 1: no interior row exceeds prominence $0.10$ and
every argmax is pinned at the port column rather than advancing.

\subsection{Wire homogeneity}

A write-count census (min/median/max writes per wire, body only) asks whether
any threshold separates the band from the carriers.

\begin{center}
\begin{tabular}{lrrl}
\toprule
build & carriers & band & separable \\
\midrule
$[3,3]$ two-carrier, roll 1 & 218/440/715 & 272/433/595 & no \\
$[1,2,3,3]$ single, band 56, roll 1 & 452/711/1121 & 449/715/1026 & no \\
$[1,2,3,3]$ single, band 256, roll 1 & 185/452/847 & 180/428/848 & no \\
$[1,2,2,3]$ single, band 256, roll 1 & 210/449/1030 & 129/426/944 & no \\
\midrule
single, band 14, \textbf{roll 0} ($n{=}16$) & 154/236/375 & 10/17/21 & \textbf{YES} \\
two-carrier, band 12, roll 0 ($n{=}16$) & 80/142/216 & 12/19/19 & \textbf{YES} \\
\bottomrule
\end{tabular}
\end{center}

Two things follow. First, \textbf{rolling matters more in the single-carrier
geometry, not less}: with one carrier per value each carrier absorbs every fold
write, so an unwritten band wire stands out against a sharper contrast, and the
roll-0 separation is cleaner than in the paired build. Rolls erase it at about
$+17\%$ gates.

Second, band 256 makes the two populations equal in size with write
distributions sitting on top of each other, so no \emph{write-count} threshold
separates them --- qualitatively better than rolling a 46-wire minority, which
a windowed census can still narrow down.

\textbf{What this does NOT buy, and an argument to retract.} An earlier draft
of this section also claimed band 256 raises the restricted degree-3 attack
from $\binom{46}{3}\approx 1.5\cdot10^4$ regressors to
$\binom{256}{3}\approx 2.8\cdot10^6$. That is arithmetically true and
strategically irrelevant: it prices an adversary that \emph{enumerates}
candidate band subsets, and \S\ref{sec:persist} exhibits one that
\emph{recovers} the band instead. Against that adversary, band width is not a
defense at any size --- it only moves where the elbow sits.

\subsection{Are the source bits good sources?}

\code{source\_stats} measures the ensemble the masks are built from: marginal
bias, pairwise correlation, the degree-3 firing rate against the $1/8$ it is
assumed to have, identically-zero terms over all offset assignments, and an
exhaustive scan for exact linear relations. Band wires and carriers both come
out well-behaved marginally (no wire above the 3-sigma floor; triple-AND means
$0.1253$ and $0.1248$ against $0.125$). The instructive difference is
structural:

\begin{center}
\begin{tabular}{lrr}
\toprule
exact wire-pair linear relations, by depth & banded build & distributed build \\
\midrule
10\,\% & --- & 21 \\
25\,\% & 1 & 23 \\
50\,\% & 1 & 2 \\
75\,\% & 1 & 0 \\
\bottomrule
\end{tabular}
\end{center}

Carriers are \emph{not} generic bits: a value that is identically $0$ on the
zero slice has two \emph{equal} shares. Identically-zero mask terms occur only
in distributed builds and only early. This is why the distributed draw now
consults a live simulation.

\section{The tradeoff: $[1,2,3,3]$ versus $[1,2,2,3]$}
\label{sec:tradeoff}

Both plans have the same atom count, hence the same fold cost, and at $n{=}128$
they build to within $340$ gates of each other on identical width. Their
degree-1 ridges are identical and neither census separates. They differ on two
axes that pull in opposite directions.

\textbf{$[1,2,2,3]$ wins the degree-1 statistic}, exactly as \eqref{eq:piling}
predicts: $0.583$ median and $0.600$ peak against $0.618$ and $0.652$, and
against the shipped build's $0.719$ and $0.749$.

\textbf{$[1,2,3,3]$ wins the degree wall.} Degree-2 terms are worth precisely
nothing to a degree-2-capable adversary --- they lie inside its span --- so
against that adversary $[1,2,2,3]$ is effectively $[1,3]$ and $[1,2,3,3]$ is
$[1,3,3]$. Measured with a full-band degree-2 predictor, rolls off so the
adversary can actually reach the band, and with the mandatory capability
control:

\begin{center}
\begin{tabular}{lrrl}
\toprule
build & ridge & $\rho$ & verdict \\
\midrule
control: degree-2 masks only & \textbf{0.3750} & 0.999 & \textbf{ALIVE} --- the adversary is capable \\
$[1,2,3,3]$ single & \textbf{0.0312} & 0.727 & dead \\
$[1,2,2,3]$ single & \textbf{0.0625} & 0.843 & twice the residual \\
\bottomrule
\end{tabular}
\end{center}

This is the ``a single degree-3 term leaves a boundary residual, two do not''
effect from the original validation, reproduced for the single-carrier decode.

\textbf{Recommendation: $[1,2,3,3]$.} It beats the shipped decode on every axis
at once --- fewer gates, fewer wires, degree-1 and degree-2 ridges both dead
against a capable adversary, and materially better statistics --- while keeping
two degree-3 terms. $[1,2,2,3]$ buys roughly $0.035$ of degree-1 margin and
spends half the degree-2 margin for it; the wall is the thing meant to hold as
adversary budgets grow. $[1,2,2,3]$ remains a one-flag alternative if the
threat model ever narrows to degree-1 statistical adversaries specifically.

\section{The band is recoverable by lifetime, and what fixes it}
\label{sec:persist}

\subsection{The attack}

Every instrument above is either per-snapshot (the two heatmaps) or a
write/read census. None of them sees the invariant that \emph{defines} the
band: a band value is a \textbf{frozen function of the input}, unchanged from
the fill to the strip. That is not incidental --- it is what the fold requires,
since a registered mask must denote the same bit at every use.

Fingerprint each wire's Boolean function at checkpoints (bit-sliced over chosen
inputs, complements identified), and match fingerprints \emph{across} wires.
Rolling relocates a variable but preserves its function, so the fingerprint
follows it for free, and the band separates from the carriers by lifetime
alone --- without ever guessing a wire set. \code{src/bin/persistence\_census}
implements the statistic: coverage of a fingerprint is the fraction of
checkpoints at which it appears on some wire.

\begin{center}
\begin{tabular}{lrrrr}
\toprule
build & band & elbow rank & coverage drop & recovered ($\ge 0.50$) \\
\midrule
$[3,3]$ two-carrier + roll 1 & 46 & \textbf{45} & $0.875 \to 0.175$ & 45 \\
$[1,2,3,3]$ single, band 256, roll 1 & 256 & \textbf{254} & $0.850 \to 0.275$ & 254 \\
\bottomrule
\end{tabular}
\end{center}

Both with rolling on, at $n{=}128$. The elbow lands on the band size, and it is
reported to survive phase-B growth. Once the population is known the rest is
cheap, and cheaper than a monomial search suggests: the carrier/band
\emph{partition} means a gate whose controls are all band wires is a mask
injection or strip, and one mixing band and carrier controls is a fold
fragment. The invariant recorded above as a safety property is a classifier.

\subsection{Retire-and-refill epochs}

The fix has to change what the attack keys on, so it must change the band
values' \emph{functions}, not their locations. One event retires a band
variable --- re-sourcing every live mask that names it, so nothing reads it ---
and then rewrites the wire; round-robin, so the population turns over every
$\text{epoch}\times\text{band}$ gaps. Ordering is the correctness argument:
release first, rewrite second, so the ledger never holds a reference across
the rewrite.

Measured at $n{=}16$, band 14, rolling on:

\begin{center}
\begin{tabular}{lrrr}
\toprule
epoch & gates & elbow drop & recovered \\
\midrule
off & 4\,698 & 0.550 & \textbf{14 / 14} \\
8 & 5\,226 & 0.075 & 14 \\
4 & 5\,719 & 0.050 & 1 \\
3 & 6\,043 & 0.050 & \textbf{0} \\
2 & 6\,786 & 0.025 & \textbf{0} \\
\bottomrule
\end{tabular}
\end{center}

The signal dies at 3--4 full band turnovers, a scale-free criterion: a variable
refilled $k$ times covers about $1/k$ of the checkpoints, so $k \gtrsim 3$ puts
every function below any useful threshold. At production ($7\,920$ source
gates, band 256) that is $\text{turnovers} = 30.9/\text{epoch}$, so epoch 10.
Confirmed:

\begin{center}
\begin{tabular}{lrr}
\toprule
$n{=}128$, band 256 & baseline & epoch 10 \\
\midrule
gates & 229\,712 & 247\,929 \ ($+7.9\%$) \\
fingerprints $\ge 0.50$ coverage & \textbf{254} & \textbf{0} \\
elbow drop & 0.575 & \textbf{0.025} \\
degree-1 ridge (median) & 0.0000 & 0.0000 \\
agreement (peak) & 0.618 (0.652) & 0.623 (0.649) \\
\bottomrule
\end{tabular}
\end{center}

So the defense is close to free on every other axis: the ridge does not move
and the statistical cost is within noise. The per-event cost is the slots
naming the retired variable, $n\,k_{\mathrm{total}}\deg/\text{band}$, so the
cost of $N$ turnovers is independent of band width --- which is why this is
$+8\%$ at production against $+29\%$ at $n{=}16$.

\textbf{Refill composition} (\code{--prod-refill-data}) barely moves the elbow:
band-internal and carrier-injecting refills are indistinguishable on it. The
argument for a nonzero rate is not statistical but structural --- a
band-internal-only recursion keeps every future band value inside the algebra
generated by the initial band --- and the counter-argument is computational:
closure is not recoverability, and following that recursion through phase A/B
is a much harder problem than fingerprinting through it. The parameter is
therefore left open rather than defaulted on this evidence.

\subsection{A three-way tension}

Trying to combine everything is what surfaced it. The reserved pivot block
(\S\ref{sec:limits}) removes every pivot from both the linear pool and the
product-source draw, so at $\text{band} = n$ there is no legal non-pivot data
wire and the fill cannot terminate; the guard now requires roughly
$\text{band} \le \tfrac34 n$. Hence:

\begin{center}
\begin{tabular}{lll}
\toprule
property & wants & compatible with epochs? \\
\midrule
write-census homogeneity (1:1 split) & band $= n$ & yes \\
provable joint uniformity (pivot block) & band $\lesssim \tfrac34 n$ & yes \\
persistence defense (epochs) & any band & --- \\
\bottomrule
\end{tabular}
\end{center}

The first two are mutually exclusive at band 256. Band 128 buys both at a
smaller band; band 256 buys homogeneity and leaves the statistics measured
rather than proved. Which to prefer depends on which adversary is weighted,
and this document does not settle it.

\section{Distributed sourcing: a negative result}
\label{sec:dist}

For completeness, since the machinery is in the tree behind
\code{--prod-src-dist}. Sourcing masks on ordinary carriers, with a write
barrier replacing the freeze, is exact and costs no extra wires --- and is
\textbf{not recommended}. At $n{=}128$ it costs $+63\%$ gates (403\,556 against
247\,326; the bill is 79\,893 forced migrations against 7\,919 ordinary
re-sources) \emph{and} leaves a faint but monotonically progress-aligned ridge
(interior median $0.0312$, $\rho = 0.995$, perm-$z$ $9.94$) where the band build
leaves none.

Eight interventions failed to move it: migration-traffic volume (a band build
with churn raised to match stays dead), churn rate (a plateau at 64$\times$
churn), source location (restricting to the sandwich's keyed-junk half changes
nothing), carrier routing, mask offsets, fill quality, dead terms, and
duplicate terms. Dumping the fitted relations shows what the leak \emph{is}:
every leaking relation has support exactly 2, the interior count is 3\,891
against the band build's 2, and four of five sampled pairs are carriers of
\emph{different} values --- a cross-value linear relation, not a value that lost
its mask (a census of bare values reads $0/256$ at every position). The
micro-mechanism is not established.

\section{Digestibility: making the fold reachable by the store}
\label{sec:digest}

The gadgetiser's output is not the deliverable; it is the input to phase~A,
whose only re-encoding channel is the frozen store, and that channel is capped
at two controls. A fold the store cannot reach is a fold that phase~A cannot
mix, so the encoding's width profile is a security parameter and not an
implementation detail.

\subsection{Narrowness is not digestibility}

\code{--prod-g57-narrow} re-realises every width-$\leq$2 fold fragment through
the g57 form, at a cost of one or two gates instead of one, with the residual
constant going to the ledger exactly as the narrow path already does. At
$n{=}128$, band 256, plan $[1,2,3,3]$, the store's match rate on the result
rises from $0.3506$ to $0.4151$ for $+4.5\%$ gates.

The rate is measured with the pure sampler --- \code{--db-dry-run} with
\code{--p-db 1.0}, every other move weight zero, \code{--db-ctrl-cap 2} --- so
the circuit never changes under the instrument and the number is a property of
the encoding rather than of a mixing trajectory.

\textbf{The mechanism is open, and an earlier explanation here was wrong.}
That explanation ran: over the wires $\{h,x,y\}$ the $X$-free span of the
g57/\textsc{cnot} vocabulary is $\langle x,\; y,\; 1 \oplus xy \rangle$, so a
width-2 conjunction with $\text{comp}=0$ is outside it and therefore invisible
to a store built from g57 circuits. The span identity is correct and worth
keeping: writing $f$ in ANF, the span is exactly
$\{f : \text{const}(f) = \text{coeff}_{xy}(f)\}$, which places $\lnot x \lnot y$
inside and $xy$, $\lnot x\, y$, $x \lnot y$ and the constant $1$ outside --- and
that is precisely why three of the four polarity cases in \code{emit\_g57\_form}
owe a deferred ledger constant while $\lnot x \lnot y$ owes none.

The \emph{conclusion} does not follow. The span describes only circuits in
which $x$ and $y$ are never written. Lift that restriction and g57 with
\textsc{cnot} generates the \emph{entire} symmetric group on three wires (all
$40\,320$ permutations, verified by closure): $h \mathrel{\oplus}= xy$ is
reachable in five gates, and even a bare \textsc{not} in five. fmix's
replacement windows are $[2,5]$ gates besides, so a lone gate is never matched
on its own. The lever is kept for the measurement; why the measurement moves is
not established.

\subsection{A width ceiling, and where it stops paying}

\code{--prod-ladder-cap $U$} ladders every fold fragment of width in $(2, U]$
down to two controls over borrowed dirty carriers, leaving wider fragments as
single gates. The same ceiling applies to slot emission: a degree-3 mask term
is a three-control gate at every inject, re-source and strip, and excluding
those would leave a residue that is \emph{attributable} --- the surviving wide
gates would be exactly the slot emissions, which name a value's mask sources
directly.

\begin{center}
\begin{tabular}{rrrrrrr}
\toprule
$U$ & gates & $\times$base & wide ($>2$) & frac & matched & rate\\
\midrule
0  & 240\,016 & 1.00 & 134\,785 & 0.562 & 151\,459 & 0.414\\
\textbf{3}  & \textbf{404\,249} & \textbf{1.68} & \textbf{93\,618} & \textbf{0.232} & \textbf{182\,136} & 0.393\\
4  & 694\,579 & 2.89 & 57\,350 & 0.083 & 139\,273 & 0.291\\
5  & 1\,041\,443 & 4.34 & 28\,635 & 0.028 & 118\,277 & 0.247\\
64 & 1\,488\,911 & 6.20 & 853 & 0.0006 & 105\,805 & 0.221\\
\bottomrule
\end{tabular}
\end{center}

Eliminating the wide gates is achievable --- from $56\%$ of the circuit to
$0.06\%$ --- and the sweep says not to buy it. \textbf{Laddering makes a gate
eligible without making it matchable at the same rate.} The match rate falls
monotonically with the ceiling, and the absolute number of successful rewrites
per unit of sampling effort peaks at $U{=}3$, some $20\%$ above the unladdered
build, before falling below it. The recommended operating point is therefore
$U{=}3$: a third of the wide gates removed and a fifth more effective store
reach, for $1.68\times$ the gates.

The $853$ width-3 gates that survive every ceiling are the
\code{--slice-zero-ccnot} preblock, three-control by design --- they are what
makes the off-slice disturbance quadratic in $x$ rather than affine. They are
not fold material.

\subsection{Diversifying CG}

A fold block emits exactly $(1 + k_{\text{total}})^{\text{arity}}$ fragments,
which is 16 for $[1,2,3,3]$ at arity 2. That is a block boundary an adversary
can find by counting and an arity it can read off the count.
\code{--prod-cg-jitter} gives a percentage of values one extra mask
term, so the count becomes $(1 + k_a)(1 + k_b) \in \{16, 20, 25\}$ and varies
with which values the source gate happens to read. Jitter adds terms and never
removes them: the operating point a build commits to is its weakest value, so
removing a term anywhere would move the commitment, while adding one only raises
a value above the floor. At $50\%$ the block-count distribution is
$\{\tfrac14, \tfrac12, \tfrac14\}$, which maximises its entropy.

The extra term is the \emph{low}-degree one. Both choices break the count
identically, but a degree-3 extra widens the fold's fragments, and that is
expensive in exactly the currency the width ceiling is spending: at $n{=}128$,
$U{=}3$, a high-degree jitter of $50\%$ cost $+14\%$ gates and $+32\%$ more wide
gates, against $+15\%$ gates and $+25\%$ wide for the low-degree version. Jitter
does not move the store's reach in either direction (matched $182\,136$ against
$182\,108$) --- it is a pure diversification spend, and $25\%$ halves the bill
if the block-count entropy is worth less than the gates.

A per-fragment ESOP re-cover was built for the same job and is \textbf{refuted}.
The identity
\[
  \operatorname{conj}(L \cup \{u\}) \;\oplus\; \operatorname{conj}(L \cup \{\lnot u\}) \;=\; \operatorname{conj}(L)
\]
is exact and does vary the block size, but the two halves share every literal
but one polarity and write the same carrier, so the twin is greppable within
the block \emph{whatever order they are emitted in}; shuffling only makes
adjacency less likely, and the test written to defend the lever found adjacent
twins anyway. Over $\mathrm{GF}(2)$ every two-term ESOP cover of a product term
\emph{is} a polarity split, so no variant of the lever escapes this. It is the
same failure mode that ruled out decoy fragments: a planted opposite-polarity
motif is worse than no lever at all.

\section{What is not proved}
\label{sec:limits}

\textbf{The piling-up numbers are measurements, not theorems.} Equation
\eqref{eq:piling} assumes the mask's source bits are \emph{jointly} uniform and
that a value's terms are variable-disjoint. \emph{The disjointness half is now
enforced in code} (\code{d9916bc8}): \code{draw\_slot} excludes every variable a
value's other live slots already name, and the retire-and-refill path
(\code{inject\_avoiding\_var}) does the same, both relaxing only when the band is
too narrow to honour it. Enforcing it is free and measurably helps --- at
$n{=}64$, paired over three gadget seeds, disjoint draws beat the previous
independent-per-term draws in $6/6$ pairs (mean $-0.010$ at $A{=}4$, $-0.006$ at
$A{=}6$), which is the right sign and order: two terms sharing a source collapse
as $w_a w_b \xor w_a w_c = w_a(w_b \xor w_c)$, paying two atoms of fold cost for
one atom of hiding.

What remains open is the \emph{joint uniformity} half. The fill establishes only
\emph{marginal} balance, one wire at a time, and the accompanying test checks
exactly that marginal --- while a mask multiplies three band wires together. Two
concrete gaps: the nonlinear fill draws each wire's pivot \emph{with} replacement
across band wires, and each wire's exclusions cover only its \emph{own} pivot, so
one wire's pivot can re-enter another's linear part. A reserved pivot block
(distinct pivots, the whole pivot set excluded from every wire's non-pivot
material) makes the band exactly uniform on $\{0,1\}^b$ for \emph{any} nonlinear
part, at zero gate cost --- see \code{docs/BAND\_INDEPENDENCE}. It is available
as \code{fill\_pivots}, and \textbf{the production preset sets it to $0$}, so
this hypothesis is not merely unproved in the shipped build, it is switched off.
It is now the \emph{only} unmet hypothesis of \eqref{eq:piling}.

Do not read the measured-versus-theory residuals as a clean estimate of that
gap: the deviation has no reliable sign, because two mechanisms push opposite
ways. \code{hmap\_stat} searches only single wires and pairwise XORs and reports
a median over interior rows, so it \emph{under}-reads a true affine adversary
(at $n{=}128$, $[1,2,3,3]$ measures $0.618$ against a theory value of $0.641$);
but each cell is a max over columns, which inflates upward, and at $n{=}64$ the
disjoint builds land \emph{above} theory ($+0.008$ at $A{=}4$, $+0.020$ at
$A{=}6$). Attributing either residual to source non-uniformity alone is not
supported. Treat \eqref{eq:piling} as a model the measurements bracket rather
than track, and note that the reduction behind it (predicting $v$ reduces to
approximating the mask) is a \textbf{lower} bound on adversary success --- an
equality only under jointly uniform sources, which is exactly what is off.

\textbf{The agreement is progress-\emph{aligned}, which ``dead at degree 1'' does
not say.} The ridge columns above are the exact-span measure, which snaps every
inexact fit to $0.5$; the agreement columns report a level but not \emph{where}
it sits. Measured with \code{stress\_battery} on the $[1,2,3,3]$ band-256 build
of \S\ref{sec:prod} (the same file: $229\,712$ gates, $512$ wires), a degree-1
\emph{bias} adversary recovers a diagonal that is both source-specific and
monotonically advancing: prominence $0.056$ against an overfit floor of $0.013$,
$\rho = 0.997$, perm-$z = 6.58$, and a decoy-source contrast of $+0.041$ (an
unrelated $C'$ reads at the floor). So the exact-span diagonal is genuinely dead
while a cheap bias adversary still reads computational progress. Single-carrier
\emph{attenuates} this relative to two-carrier ($0.135 \to 0.056$ prominence) but
does not remove it, and $6$M \code{fmix} transport moves with a $2.4\times$ growth
leave it structurally intact ($\rho\ 0.996 \to 0.997$, contrast $+0.081 \to
+0.084$): it is a value-preserving leak, and every mixing move preserves values.
See \code{docs/STRESS\_BATTERY}.

\textbf{One-probe independence is lost in single-carrier mode.} A two-carrier
share is uniform \emph{and} independent of the value, so a single probe learns
nothing; a lone carrier holds $v \xor \text{masks} \xor \kappa$, correlated with
the value at the piling-up rate. The ledger constant does not rescue this: it is
a compile-time bit, identical for every input, so it flips the \emph{sign} of
that correlation rather than its magnitude, and one bit is learnable from a
handful of samples. This is better against pair-searching and affine
adversaries --- which every heatmap here is --- and worse against a strict
one-probe adversary.

\textbf{Agreement is not security.} $0.62$ agreement on a bit is a $0.12$ bias,
amplifiable across snapshots and across correlated cells. These are lower
bounds from a bounded predictor family, and the construction has no upper bound
of any kind at present.

\textbf{Scope.} Everything measured here is execution-based, per-instance, and
single-snapshot: one $(C,G)$ pair, one time-pair per cell, uniform inputs with
the slice pinned. It says nothing about the distributional game $G(C_1)$ versus
$G(C_2)$, about joint multi-snapshot leaks, about subcube-conditional
adversaries, or about search/SAT-guided attacks. The static/syntactic channel is
covered only by the ad-hoc censuses reported above.

\section{Options}

\begin{center}
\begin{tabular}{llp{0.46\textwidth}}
\toprule
flag & env & meaning \\
\midrule
\code{--prod-k} / \code{--prod-deg} & \code{PROD\_K} / \code{PROD\_DEG} & base terms and their degree \\
\code{--prod-k-hi} / \code{--prod-deg-hi} & \code{PROD\_K\_HI} / \code{PROD\_DEG\_HI} & tower terms and their degree \\
\code{--prod-single} & \code{PROD\_SINGLE} & single-carrier decode (one linear term) \\
\code{--prod-band} & \code{PROD\_BAND} & band width; 0 = match the value count \\
\code{--prod-rsrc} & \code{PROD\_RSRC} & re-source moves per gap; load-bearing under \code{--prod-single} \\
\code{--prod-roll} & \code{PROD\_ROLL} & band relocations per gap \\
\code{--prod-fill-nl} & \code{PROD\_FILL\_NL} & nonlinear cascade fill \\
\code{--prod-max-width} & \code{PROD\_MAX\_WIDTH} & narrow mode (not recommended) \\
\code{--prod-epoch} & \code{PROD\_EPOCH} & gaps between retire-and-refill events; kills the persistence signal (\S\ref{sec:persist}) \\
\code{--prod-refill-data} & \code{PROD\_REFILL\_DATA} & \% of refill sources drawn from carriers rather than band wires \\
\code{--prod-fill-pivots} & \code{PROD\_FILL\_PIVOTS} & reserved pivot block; needs band $\lesssim \tfrac34 n$ \\
\code{--prod-src-dist} & \code{PROD\_SRC\_DIST} & distributed sourcing (not recommended, \S\ref{sec:dist}) \\
\bottomrule
\end{tabular}
\end{center}

The plan is the multiset built as $k$ copies of $\deg$ then $k_{hi}$ of
$\deg_{hi}$; the encoding is enabled iff $k>0$ or $k_{hi}>0$ (the degree flags
alone are inert), and \code{--prod-deg 1} is silently clamped to 2. The two
recommended settings are

\begin{center}
\begin{tabular}{ll}
\toprule
$[1,2,3,3]$ (recommended) & \code{--prod-single 1 --prod-k 1 --prod-deg 2 --prod-k-hi 2 --prod-deg-hi 3} \\
$[1,2,2,3]$ & \code{--prod-single 1 --prod-k 2 --prod-deg 2 --prod-k-hi 1 --prod-deg-hi 3} \\
\bottomrule
\end{tabular}
\end{center}

both with \code{--prod-fill-nl 2 --prod-roll 1 --prod-g57-narrow 1
--prod-ladder-cap 3 --prod-cg-jitter 50 --prod-epoch 5 --prod-refill-data 50},
which is \code{ProdConfig::production\_single} in the tree.

\textbf{That preset is now the DEFAULT, and this is a change.} Every
\code{--prod-*} lever used to default to $0$ --- and the preset itself sat in
the tree with no callers, so the levers this document calls "the validated
production setting" were off in every circuit anyone actually built. Both entry
points (\code{sss} and \code{gen\_sandwich\_gadget}) now construct their config
\emph{from} the preset and let an explicitly passed flag or environment
variable override one field, so the values live in exactly one place and cannot
drift out of the build again. A unit test pins them.

Two consequences worth stating plainly. \code{--prod-k 0} is now the way to ask
for no encoding at all; and \code{--prod-band 0} no longer means the old auto
rule $\max(\lceil\sqrt{4\,n\,k_{\mathrm{total}}}\,\rceil, 6, \max\deg+3)$
($56$ at the production sandwich) but \emph{match the value count}, which is
the $1{:}1$ split the write census needs. Here $n$ is the \emph{entry point's}
value count --- the sandwich width $2n$, not the source circuit's $n$ --- so the
production build is $256$ carriers and $256$ band wires. Pass a number to
\code{--prod-band} for any other sizing. \code{--prod-single 0} still
reproduces the two-carrier build.

\end{document}
